\section{Ubicación del curso y pregunta central}

\subsection{Por qué importa esta clase}
Esta clase no busca presentar otro «modelo más grande», sino centrar la atención en un problema más de ingeniería: cuando el modelo entra a un entorno web real, cómo pasar de «saber responder preguntas» a «saber completar tareas». Ruslan resume el tema en tres puntos: evaluación, búsqueda y escalamiento de datos. Estos tres puntos son justamente donde la mayoría de los sistemas de Agentes encuentran primero obstáculos al llevarlos a producción.

\begin{importantbox}{Hilo conductor de la clase}
\begin{enumerate}
    \item \textbf{VisualWebArena}: cómo evaluar Agentes web multimodales en un entorno reproducible;
    \item \textbf{Inference-Time Search}: cómo lograr que el Agente sea mejor para probar, equivocarse y retroceder durante la inferencia;
    \item \textbf{Data Scaling}: cómo seguir mejorando la capacidad del Agente en condiciones de escasez de datos.
\end{enumerate}
\end{importantbox}

\begin{knowledgebox}{El cambio de paradigma de «conversar» a «ejecutar»}
Cuando la tarea pasa de responder preguntas a operar una página web, el modelo debe procesar al mismo tiempo la entrada visual, los cambios de estado, las restricciones de acción, los objetivos de largo plazo y la recuperación ante fallas. Este problema se acerca más a «Sequential Decision Making» que a la generación de texto en un solo turno.
\end{knowledgebox}

\begin{figure}[H]
\centering
\includegraphics[width=0.82\textwidth]{frame-01-visualwebarena.jpg}
\caption{Video aproximadamente en 00:06:50: el ponente presenta VisualWebArena y enfatiza la evaluación de Agentes multimodales en un entorno de interacción real}
\end{figure}

\subsection{Qué aporta esta clase directamente a quienes trabajan en la práctica}
El ponente insiste en un punto: el cuello de botella de los sistemas de Agentes no está solo en los parámetros del modelo, sino en la capacidad sistémica. Incluso si el modelo base es más fuerte, sin búsqueda, verificadores y datos de interacción de alta calidad, el límite superior de desempeño se alcanza muy pronto. A la inversa, aunque el modelo base no sea el más fuerte, con solo acertar en estas tres líneas es posible obtener beneficios considerables.

\begin{warningbox}{Error frecuente}
Equiparar directamente «la mejora de la capacidad del modelo» con «el aumento de la tasa de finalización de tareas» suele ser un error. En las tareas web, el espacio de estados es grande, los pasos son largos y el costo de los errores intermedios es alto, por lo que la mejora de la capacidad en un solo paso no se traduce automáticamente en una mayor tasa de éxito en cadenas largas.
\end{warningbox}

\subsection{Resumen de la sección}
Esta clase se centra en los tres pilares de la ingeniería de Agentes: evaluación, búsqueda y datos. No le interesa la calidad de una respuesta aislada, sino la tasa de éxito y la escalabilidad de las tareas de cadena larga.

